\documentclass[12pt]{article}
\usepackage{notes-common}
\usetikzlibrary{calc}
\geometry{margin=0.35in}
\AtBeginDocument{\fontsize{14}{17}\selectfont}

% Complete / filled-in reference version. No skeleton (blank) companion
% was requested for this one -- key only, per Seth 2026-10-19.
% Fleshed out from "Circuits review - Rough.txt": charge on a capacitor
% at an instant during charging (incl. a parallel-capacitor variant with
% the knowns/unknowns swapped), then the loop rule at steady state, two
% structurally different circuits. All four answers verified with a
% script before typing in. Numbers reused from circuits-review-practice.tex
% (in-class-practice), which was built from the same rough notes first.
\newcommand{\ANS}[1]{\textcolor{accent}{#1}}

% ---- same schematic pieces as circuits-devices-notes-key, plus a
% capacitor symbol (not needed by any earlier notes file) ----
\newcommand{\battv}[1]{% vertical battery symbol centered at (#1)
  \draw[line width=0.7pt] ($(#1)+(0,0.55)$) -- ($(#1)+(0,0.12)$);
  \draw[line width=0.7pt] ($(#1)+(0,-0.12)$) -- ($(#1)+(0,-0.55)$);
  \draw[line width=1.1pt] ($(#1)+(-0.35,0.12)$) -- ($(#1)+(0.35,0.12)$);
  \draw[line width=2.4pt] ($(#1)+(-0.2,-0.12)$) -- ($(#1)+(0.2,-0.12)$);
  \node[scale=0.75] at ($(#1)+(0.6,0.12)$) {$+$};
  \node[scale=0.75] at ($(#1)+(0.6,-0.12)$) {$-$};
}
\newcommand{\reszig}[1]{% horizontal resistor zigzag centered at (#1)
  \fill[white] ($(#1)+(-0.65,-0.32)$) rectangle ($(#1)+(0.65,0.32)$);
  \draw ($(#1)+(-0.6,0)$)
    -- ($(#1)+(-0.42,0.22)$) -- ($(#1)+(-0.24,-0.22)$)
    -- ($(#1)+(-0.06,0.22)$) -- ($(#1)+(0.12,-0.22)$)
    -- ($(#1)+(0.3,0.22)$) -- ($(#1)+(0.48,-0.22)$)
    -- ($(#1)+(0.6,0)$);
}
\newcommand{\capsym}[1]{% horizontal capacitor (two plates), centered at (#1)
  \fill[white] ($(#1)+(-0.35,-0.4)$) rectangle ($(#1)+(0.35,0.4)$);
  \draw[line width=0.7pt] ($(#1)+(-0.35,0)$) -- ($(#1)+(-0.1,0)$);
  \draw[line width=0.7pt] ($(#1)+(0.1,0)$) -- ($(#1)+(0.35,0)$);
  \draw[line width=1.8pt] ($(#1)+(-0.1,-0.3)$) -- ($(#1)+(-0.1,0.3)$);
  \draw[line width=1.8pt] ($(#1)+(0.1,-0.3)$) -- ($(#1)+(0.1,0.3)$);
}
\newcommand{\capsymv}[1]{% vertical capacitor, centered at (#1)
  \fill[white] ($(#1)+(-0.4,-0.35)$) rectangle ($(#1)+(0.4,0.35)$);
  \draw[line width=0.7pt] ($(#1)+(0,-0.35)$) -- ($(#1)+(0,-0.1)$);
  \draw[line width=0.7pt] ($(#1)+(0,0.1)$) -- ($(#1)+(0,0.35)$);
  \draw[line width=1.8pt] ($(#1)+(-0.3,-0.1)$) -- ($(#1)+(0.3,-0.1)$);
  \draw[line width=1.8pt] ($(#1)+(-0.3,0.1)$) -- ($(#1)+(0.3,0.1)$);
}

\begin{document}

\noteshead{Circuits Review: Capacitors \& Kirchhoff's Rules \textcolor{accent}{(completed)}}{}

\goal{
\begin{itemize}[itemsep=5pt,parsep=2pt,topsep=2pt]
\item See that the loop rule applies to a capacitor exactly like it does to
a resistor or a battery --- it just contributes a $Q/C$ term.
\item Use that at \emph{some instant during charging} (Problems 1--2), where
the general relation $\varepsilon = IR + Q/C$ holds at every moment, not
just $t=0$ or $t=\infty$.
\item Use it again at \emph{steady state} (Problems 3--4), where a fully
charged capacitor carries no current --- and see that this sometimes makes
a series resistor's value completely irrelevant, and sometimes doesn't.
\end{itemize}}

%============================================================
\nsec{Problem 1 --- charge at an instant}

A capacitor $C$ is charging through a resistor $R$, connected to a battery
$\varepsilon$, all in a single loop. At a certain moment, while charge is
still building up, the current through the resistor is $I$.

\begin{center}
\begin{tikzpicture}[x=1cm,y=1cm]
  \draw (0,0) -- (0,2) -- (5,2) -- (5,0) -- cycle;
  \battv{0,1} \node[scale=0.8, anchor=east] at (-0.55,1) {$\varepsilon$};
  \reszig{2.5,2} \node[scale=0.8, anchor=south] at (2.5,2.45) {$R$};
  \capsymv{5,1} \node[scale=0.8, anchor=west] at (5.55,1) {$C$};
\end{tikzpicture}
\end{center}

Write the loop rule for this circuit (valid at \emph{any} instant, charging
or not) and solve for the charge on the capacitor at this moment.

{\small
The loop rule, one resistor and one capacitor:
\[
\varepsilon - IR - \dfrac{Q}{C} = 0.
\]}

\[
\boxed{Q = C(\varepsilon - IR)}
\]

\textbf{Now with} $\varepsilon=9.0$ V, $R=2.0\ \text{k}\Omega$,
$C=4.0\ \mu\text{F}$, $I=1.5$ mA:

\[
Q = (4.0\ \mu\text{F})\big[9.0\text{ V} - (1.5\text{ mA})(2.0\ \text{k}\Omega)\big]
= (4.0\ \mu\text{F})(6.0\text{ V}) = \ANS{24\ \mu\text{C}}
\]

\checkpoint{This current, $I=1.5$ mA, is less than the circuit's very first
current, $I_0=\varepsilon/R$. What is $I_0$ here, and why must the given
$I$ always be smaller than it?}

{\small\color{accent} $I_0=\varepsilon/R=9.0\text{ V}/2.0\ \text{k}\Omega=
\ANS{4.5\text{ mA}}$. The charging current only ever \emph{decays} from
$I_0$ (as the capacitor fills and starts opposing it) --- it never starts
small and grows. So any instant after $t=0$ must have $I<I_0$, as it does
here.}

%============================================================
\nsec{Problem 2 --- same idea, with two capacitors in parallel}

Two capacitors, $C_1=4.0\ \mu\text{F}$ and $C_2=6.0\ \mu\text{F}$, are
wired in parallel with each other. That parallel pair is in series with a
resistor $R=2.0\ \text{k}\Omega$ and a battery $\varepsilon=15$ V.

\begin{center}
\begin{tikzpicture}[x=1cm,y=1cm]
  \draw (0,0) -- (0,2) -- (1.3,2) -- (1.3,2.9) -- (4.7,2.9) -- (4.7,2) -- (5,2);
  \draw (1.3,2) -- (1.3,1.1) -- (4.7,1.1) -- (4.7,2);
  \draw (5,2) -- (7,2) -- (7,0) -- (0,0);
  \battv{0,1} \node[scale=0.8, anchor=east] at (-0.55,1) {$\varepsilon$};
  \capsym{3,2.9} \node[scale=0.8, anchor=south] at (3,3.35) {$C_1$};
  \capsym{3,1.1} \node[scale=0.8, anchor=north] at (3,0.65) {$C_2$};
  \reszig{6,2} \node[scale=0.8, anchor=south] at (6,2.45) {$R$};
\end{tikzpicture}
\end{center}

At a certain moment during charging, the charge on $C_1$ is measured to be
$Q_1=20\ \mu\text{C}$. Find (a) the voltage across the parallel pair at
that moment, (b) the current through $R$, and (c) the charge $Q_2$ on
$C_2$.

{\small
Parallel elements share the same voltage, so start from $C_1$ alone:
\[
\Delta V = \dfrac{Q_1}{C_1}.
\]
Then the loop rule gives the current, and $Q_2$ follows the same way $Q_1$
did:
\[
\varepsilon - IR - \Delta V = 0 \qquad Q_2 = C_2\,\Delta V.
\]}

\textbf{(a)}
\[
\Delta V = \dfrac{20\ \mu\text{C}}{4.0\ \mu\text{F}} = \ANS{5.0\text{ V}}
\]
\textbf{(b)}
\[
I = \dfrac{\varepsilon - \Delta V}{R} = \dfrac{15\text{ V}-5.0\text{ V}}{2.0\ \text{k}\Omega}
= \ANS{5.0\text{ mA}}
\]
\textbf{(c)}
\[
Q_2 = (6.0\ \mu\text{F})(5.0\text{ V}) = \ANS{30\ \mu\text{C}}
\]

\checkpoint{$C_1$ and $C_2$ have different capacitances but the \emph{same}
charge-to-voltage behavior was used for both in part (c). What justifies
using one formula, $Q=C\,\Delta V$, for each capacitor separately here,
instead of needing some combined parallel-capacitance step?}

{\small\color{accent} They're in parallel, so by definition they share the
same $\Delta V$ --- there's nothing to combine first. (If the question had
asked for the charge drawn from the \emph{battery}, that total would need
$C_{\text{p}}=C_1+C_2$; but each capacitor's \emph{own} charge only ever
needs its own $C$ and the shared $\Delta V$.)}

\clearpage
%============================================================
\nsec{Problem 3 --- steady state, resistor values don't matter}

A battery $\varepsilon=10$ V is in series with two resistors,
$R_1=1.5\ \text{k}\Omega$ and $R_2=2.5\ \text{k}\Omega$, and a capacitor
$C=3.0\ \mu\text{F}$ --- all in one single loop. The circuit has been
connected a long time (steady state).

\begin{center}
\begin{tikzpicture}[x=1cm,y=1cm]
  \draw (0,0) -- (0,2) -- (6,2) -- (6,0) -- cycle;
  \battv{0,1} \node[scale=0.8, anchor=east] at (-0.55,1) {$\varepsilon$};
  \reszig{2,2} \node[scale=0.8, anchor=south] at (2,2.45) {$R_1$};
  \reszig{4,2} \node[scale=0.8, anchor=south] at (4,2.45) {$R_2$};
  \capsymv{6,1} \node[scale=0.8, anchor=west] at (6.55,1) {$C$};
\end{tikzpicture}
\end{center}

What is the charge $Q$ on the capacitor?

{\small
At steady state the fully charged capacitor carries no current --- and
since this is a single loop, \emph{nothing else} can carry current either:
$I=0$ everywhere. Both resistors drop $IR=0$, no matter what $R_1$ and
$R_2$ actually are.
\[
\varepsilon - (0)R_1 - (0)R_2 - \dfrac{Q}{C} = 0 \quad\Longrightarrow\quad
\varepsilon = \dfrac{Q}{C}.
\]}

\[
\boxed{Q = C\varepsilon}
\]

\textbf{Now with the numbers above}:
\[
Q = (3.0\ \mu\text{F})(10\text{ V}) = \ANS{30\ \mu\text{C}}
\]

\checkpoint{Suppose $R_1$ and $R_2$ had been given totally different
values. Would $Q$ change?}

{\small\color{accent} No --- not at all. $R_1$ and $R_2$ never appear in
the boxed result; they dropped out the moment $I=0$. Any series resistor
(or chain of them) in a single loop like this one becomes completely
irrelevant to the final charge once the capacitor is fully charged. This
is the opposite situation from Problem 4, where a resistor's value very
much does matter --- the difference is whether current has somewhere else
to go.}

%============================================================
\nsec{Problem 4 --- steady state, where resistor values do matter}

A battery $\varepsilon=12$ V is in series with a resistor
$R_1=2.0\ \text{k}\Omega$. That combination feeds a parallel pair: a
second resistor $R_2=4.0\ \text{k}\Omega$ and a capacitor
$C=5.0\ \mu\text{F}$. The circuit has reached steady state.

\begin{center}
\begin{tikzpicture}[x=1cm,y=1cm]
  \draw (0,0) -- (0,2) -- (3,2);
  \draw (3,2) -- (3.3,2) -- (3.3,2.9) -- (6.7,2.9) -- (6.7,2) -- (7,2) -- (7,0) -- (0,0);
  \draw (3.3,2) -- (3.3,1.1) -- (6.7,1.1) -- (6.7,2);
  \battv{0,1} \node[scale=0.8, anchor=east] at (-0.55,1) {$\varepsilon$};
  \reszig{1.5,2} \node[scale=0.8, anchor=south] at (1.5,2.45) {$R_1$};
  \reszig{5,2.9} \node[scale=0.8, anchor=south] at (5,3.35) {$R_2$};
  \capsym{5,1.1} \node[scale=0.8, anchor=north] at (5,0.65) {$C$};
\end{tikzpicture}
\end{center}

Find (a) the current through $R_1$ (and through $R_2$), and (b) the charge
on the capacitor.

{\small
This time the capacitor being open at steady state does \emph{not} kill
all the current --- $R_2$ gives it another way around. Current flows
through the loop $\varepsilon$--$R_1$--$R_2$ instead:
\[
\varepsilon - IR_1 - IR_2 = 0 \quad\Longrightarrow\quad I=\dfrac{\varepsilon}{R_1+R_2}.
\]
The capacitor is in parallel with $R_2$, so it ends up at $R_2$'s voltage:
\[
\Delta V_C = IR_2 \qquad Q=C\,\Delta V_C.
\]}

\textbf{(a)}
\[
I = \dfrac{12\text{ V}}{2.0\ \text{k}\Omega+4.0\ \text{k}\Omega} = \ANS{2.0\text{ mA}}
\]
\textbf{(b)}
\[
\Delta V_C = (2.0\text{ mA})(4.0\ \text{k}\Omega) = 8.0\text{ V} \qquad
Q = (5.0\ \mu\text{F})(8.0\text{ V}) = \ANS{40\ \mu\text{C}}
\]

\checkpoint{In Problem 3, every resistor's value was irrelevant to $Q$.
Here, $R_1$ and $R_2$ both clearly matter. What is the one structural
difference between the two circuits that explains this?}

{\small\color{accent} In Problem 3 the capacitor is the \emph{only} path
back to the battery, so blocking it blocks the whole loop, $I=0$
everywhere, and no resistor drop survives to matter. Here $R_2$ sits
\emph{in parallel} with the capacitor, giving current a second route that
doesn't go through $C$ --- so current keeps flowing through $R_1$ and
$R_2$ even once the capacitor is fully charged, and their values set
exactly how much voltage lands across it.}

\end{document}
